\documentclass[12pt]{article}
\usepackage{amsmath}
\usepackage{amsfonts}
\usepackage{amssymb}
\usepackage{physics}
\usepackage{geometry}
\usepackage{parskip}
\geometry{a4paper, margin=0.7in}


\title{02YMECH - Solution 10}
\date{Assigned for the week of Nov 24, 2025}
\author{}

\begin{document}
\maketitle

\section*{Solutions}

\begin{enumerate}

\item At the equator, the Coriolis acceleration is
\[
\mathbf a_C = -2\,\boldsymbol{\Omega}\times\mathbf v,
\]
with \(\boldsymbol{\Omega}\) horizontal and pointing north.

The apple falls vertically with
\[
v_z(t)=-gt.
\]
Thus the Coriolis acceleration is purely horizontal (eastward) with magnitude
\[
a_C = 2\Omega g t.
\]

Horizontal velocity:
\[
v_x(t)=\int_0^t 2\Omega g \tau\,d\tau=\Omega g t^2.
\]

Horizontal displacement:
\[
x(t)=\int_0^t \Omega g \tau^2\,d\tau=\frac{1}{3}\Omega g t^3.
\]

The fall time from height \(h\) is
\[
t_f=\sqrt{\frac{2h}{g}}.
\]

Hence
\[
x=\frac{1}{3}\Omega g\left(\frac{2h}{g}\right)^{3/2}
=\frac{2\sqrt{2}}{3}\,\Omega\,\frac{h^{3/2}}{\sqrt{g}}.
\]

With \(h=4~\mathrm{m}\), \(\Omega=7.29\times10^{-5}~\mathrm{s^{-1}}\),
\[
x \approx 2\times10^{-4}\ \mathrm{m} = 0.2\ \mathrm{mm}.
\]

The deviation is negligible.  
The phrase “the apple does not fall far from the tree” remains valid :)

\newpage

\item 
\begin{enumerate}
\item[(a)] In the rotating frame, the forces along the pipe are due to the centrifugal force.
The equation of motion is
\[
m\ddot r = m\Omega^2 r.
\]

(Coriolis force is perpendicular to the pipe and constrained by the walls.)

\item[(b)] Solution:
\[
r(t)=A e^{\Omega t}+B e^{-\Omega t}.
\]
Using \(r(0)=r_0\), \(\dot r(0)=0\),
\[
r(t)=r_0\cosh(\Omega t).
\]

\item[(c)] Time to reach \(r=R\):
\[
R=r_0\cosh(\Omega t)
\;\Rightarrow\;
t=\frac{1}{\Omega}\cosh^{-1}\!\left(\frac{R}{r_0}\right).
\]

\item[(d)] Although the particle is initially at rest in the rotating frame, it experiences a centrifugal force \(m\Omega^2 r\) directed outward. This force does work on the particle, causing radial acceleration and increasing kinetic energy in the rotating frame.
\end{enumerate}

\end{enumerate}

\end{document}
